\section{Literature Review}
\label{sec:literature_review}

\subsection{Vehicle Routing Problem: Foundations and Advances}

The Vehicle Routing Problem (VRP), first introduced by \cite{Dantzig1959}, has been extensively studied in operations research. The classical Capacitated VRP (CVRP) seeks to serve a set of customers with a fleet of capacitated vehicles at minimum cost. \cite{Toth2014} provide a comprehensive review of exact and approximate solution methods.

\textbf{Exact Algorithms.} Branch-and-cut-and-price algorithms \cite{Fukasawa2006} can solve instances with up to 100 customers optimally. However, computation time grows exponentially with problem size, making exact methods impractical for large-scale humanitarian applications.

\textbf{Heuristic Algorithms.} Clarke and Wright's savings algorithm \cite{Clarke1964} remains influential for its simplicity and effectiveness. More sophisticated heuristics include:
\begin{itemize}
\item Insertion heuristics \cite{Solomon1987}
\item Sweep algorithm \cite{Gillett1974}
\item Tabu search \cite{Gendreau1994}
\item Simulated annealing \cite{Osman1993}
\item Genetic algorithms \cite{Potvin1996}
\end{itemize}

\textbf{Metaheuristics.} Recent advances focus on Adaptive Large Neighborhood Search (ALNS) \cite{Ropke2006}, which dynamically adapts search operators based on their historical performance. ALNS has demonstrated superior performance on benchmark VRP instances \cite{Pisinger2010}.

\subsection{Stochastic Vehicle Routing Problems}

Stochastic VRPs (SVRPs) extend the classical VRP by incorporating uncertainty in travel times, demands, or service times. \cite{Gendreau2014} classify SVRPs into three categories:

\textbf{SVRP with Stochastic Demands.} \cite{Dror1995} introduce chance-constrained models for demand uncertainty. \cite{Bertsimas1992} propose a priori optimization with recourse, where routes are planned in advance but can be adjusted based on realized demands.

\textbf{SVRP with Stochastic Travel Times.} \cite{Laporte1992} formulate the traveling salesman problem with stochastic travel times. \cite{Kenyon2003} develop branch-and-cut algorithms for the SVRP with time windows. Our work builds on these foundations while emphasizing reliability constraints critical for humanitarian applications.

\textbf{Dynamic SVRPs.} \cite{Pillac2013} review dynamic and stochastic VRPs where information is revealed over time. While our model adopts a static planning approach (consistent with humanitarian practice), we incorporate scenario-based stochastic programming to capture uncertainty.

\subsection{Humanitarian Logistics}

Humanitarian logistics presents unique challenges compared to commercial logistics \cite{VanWassenhove2016}. Key distinctions include:

\textbf{Objective Functions.} Unlike commercial logistics focusing on cost minimization, humanitarian logistics prioritizes \cite{Kovacs2012}:
\begin{itemize}
\item Minimizing suffering (urgency)
\item Maximizing equity (fair access)
\item Minimizing response time
\end{itemize}

\textbf{Last-Mile Distribution.} \cite{Balcik2008} study last-mile delivery in humanitarian relief chains. \cite{Pedraza-Martinez2014} examine vehicle fleet management in humanitarian logistics, identifying unique constraints in resource-constrained environments.

\textbf{Equity Considerations.} \cite{Gralla2014} develop metrics for equity in humanitarian delivery, showing that purely efficiency-focused models may leave vulnerable populations underserved. Our reliability constraints address this concern by ensuring minimum service standards for all locations.

\subsection{Road Network Reliability}

Road network reliability analysis draws from transportation engineering and infrastructure resilience literature:

\textbf{Reliability Measures.} \cite{Bell2000} define travel time reliability as the probability that a trip completes within a specified time threshold. \cite{Nicholson2006} distinguish between:
\begin{itemize}
\item Reliability: Probability of network functionality
\item Vulnerability: Consequences of network failure
\item Resilience: Ability to recover from disruptions
\end{itemize}

\textbf{Modeling Approaches.} \cite{Taylor2000} develop probit-based models for travel time variability. \cite{Lo2006} propose game-theoretic approaches for network reliability under stochastic demand. Our work applies these concepts to humanitarian routing in underdeveloped road networks.

\subsection{Multi-Depot VRP}

Multi-Depot VRP (MDVRP) extends the classical VRP to multiple starting depots. \cite{Polacek2008} review exact and heuristic approaches:

\textbf{Clustering-Based Methods.} \cite{Gillett1974} propose clustering customers before route construction. \cite{Cordeau1997} develop a parallel tabu search for MDVRP. Our approach builds on these foundations with stochastic considerations.

\textbf{Coordination Mechanisms.} \cite{Crevier2007} study MDVRP with inter-depot routes, allowing vehicles to transfer loads between depots. This is particularly relevant for humanitarian logistics where inter-warehouse collaboration can improve efficiency.

\subsection{Research Gaps}

Despite extensive literature on VRP and stochastic optimization, several gaps remain:

\textbf{Gap 1: Reliability-Driven Routing.} Most SVRP literature focuses on expected cost minimization with chance constraints. Limited work addresses reliability as a primary objective, particularly with high reliability thresholds (e.g., $\alpha \geq 0.85$) critical for humanitarian applications.

\textbf{Gap 2: Data-Driven Uncertainty Modeling.} While theoretical SVRP models exist, limited work applies these models with real-world road network data from humanitarian contexts. Our study uses Somalia road network data to fit probability distributions, bridging theory and practice.

\textbf{Gap 3: Multi-Depot Stochastic Coordination.} MDVRP literature predominantly addresses deterministic problems. Stochastic multi-depot routing with coordinated reliability constraints remains underexplored.

\textbf{Gap 4: Practical Humanitarian Implementation.} Academic models often ignore practical constraints (communication limitations, field operator training, implementation feasibility). Our work addresses these considerations through algorithm design and validation.

\subsection{Paper Positioning}

This paper contributes to the literature by:
\begin{enumerate}
\item Formulating a reliability-focused SVRP with high reliability thresholds
\item Incorporating empirical road network data from humanitarian contexts
\item Developing coordinated multi-depot strategies under uncertainty
\item Providing practical algorithms implementable in resource-constrained environments
\item Demonstrating significant improvements over deterministic baselines
\end{enumerate}

Our work extends \cite{Gendreau2014}'s SVRP framework with reliability constraints, builds on \cite{Ropke2006}'s ALNS for humanitarian applications, and contributes to the growing field of humanitarian logistics \cite{Kovacs2012}.
